\section{Four-dimensional regularity up to the boundary}\label{four:sec:regularity}

The second equation has quadratic growth in both $\grad Z$ and
$\Hess f$.  A bound for its ellipticity constants therefore does not
close the regularity argument.  We first prove an independent scalar
estimate which gives a positive Morrey exponent for $|\Hess f|^2$.
That exponent permits a measure correction in the second equation.
All constants in this section may depend arbitrarily on a fixed $N$;
none of them is subsequently used as a polynomial bound $\Pi_N$.
In coordinate calculations $D_x$, or $D$ when there is no ambiguity,
denotes ordinary differentiation, whereas $D=1+l^2|\grad f|^2$ in
the system retains its earlier meaning.

\subsection{An independent scalar estimate}

\begin{theorem}[Measurable axial coefficient]\label{four:thm:axial-regularity}
Let $U$ be a four-dimensional coordinate ball $B_2$, or a half-ball
$B_2^+=B_2\cap\{x_4>0\}$.  Assume that the metric is smooth,
$C_g^{-1}\Id\le (g_{ij})\le C_g\Id$, and that its coordinate
$C^2$ norm is at most $C_g$.  Assume also uniformly bounded $C^3$
norms for the charts and inverse charts.  On a half-ball use boundary
normal coordinates, so $g_{44}=1$ and $g_{a4}=0$ for $a<4$.
Let $z$ be $C^3$, up to the flat face when present, and constant on
that face.  Suppose, almost everywhere, that
\begin{equation}\label{four:reg:axial-equation}
 \begin{gathered}
 A_z:\Hess_g z=G,\qquad
 A_z=\Id-(1-\widetilde\chi)e_z\otimes e_z,\\
 e_z=\frac{\grad z}{|\grad z|},\qquad
 0<\chi_0\le\widetilde\chi\le1.
 \end{gathered}
\end{equation}
Here $\widetilde\chi$ and $G$ are measurable.  At $\grad z=0$,
any measurable unit axis making the equation true is allowed.
If $|\grad z|\le M_1$ and $|G|\le N_1$, then there are
$\alpha\in(0,1)$, $r_0>0$ and $C<\infty$ such that
\begin{equation}\label{four:eq:axial-estimate}
 \|\grad z\|_{C^\alpha(U_1)}\le C(M_1+N_1),\qquad
 \int_{B_r(x)\cap U}|\Hess_g z|^2\,\dd V_g
       \le C(M_1+N_1)^2r^{2+2\alpha}
\end{equation}
for $x\in\overline{U_1}$ and $0<r\le r_0$, where $U_1=B_1$
or $B_1^+$.  The constants depend only on $\chi_0$, the stated
geometry bounds and the separation from the curved patch edges.
The $C^\alpha$ norm may equivalently be taken for the coordinate
gradient components.  In particular, no modulus of continuity of
$G$ or $\widetilde\chi$, and no bound for a derivative of either,
is assumed.
\end{theorem}

\begin{proof}
We prove the assertion on uniformly small full balls, using a
reflection at the flat face.  Rescaling back and a finite local
cover give the stated form.
The argument first establishes a local Hessian estimate and a dichotomy:
on a smaller ball, either the gradient bound decreases or the solution
is close to a nonconstant affine function. A fixed-axis estimate then
improves the affine approximation. Iterating these alternatives gives
H\"older gradient control, from which the local Hessian estimate yields
the Morrey bound.

\paragraph{Doubling and a quantitative Hessian estimate.}
Subtract the constant boundary value.  With
$R_4=\operatorname{diag}(1,1,1,-1)$, extend the tangential metric
block evenly and set $\widetilde z(x)=-z(R_4x)$ below the plane.
Tangential derivatives of $z$ vanish on the face and the two normal
derivatives agree.  Thus $\widetilde z$ is $C^1$; its distributional
coordinate Hessian has no surface measure.  The doubled metric is
Lipschitz and has the prescribed smooth bounds on each closed side.
On the reflected side the gradient and covariant Hessian are
$-R_4\grad z$ and $-R_4(\Hess_g z)R_4$, respectively.  Extend
$\widetilde\chi$ evenly and $G$ oddly.  Equation
\eqref{four:reg:axial-equation} then holds almost everywhere on the double.
Each individual doubled function is locally $W^{2,p}$ for every
finite $p$, although these norms are not yet uniformly bounded.

Make a constant linear coordinate change so that the metric at the
center of a ball is the identity.  A plane in the double remains a
plane under this change.  If the ball is sufficiently small, the
coordinate principal matrix $a=(a^{ij})$ in
\eqref{four:reg:axial-equation} satisfies
\begin{equation}\label{four:reg:defect}
 |\Id-a(x)|_F\le 1-\chi_0/2=:\kappa<1.
\end{equation}
The reason for the strict gap is the exact identity
$|\Id-A_z|_F=1-\widetilde\chi$ in an orthonormal frame;
the small metric oscillation consumes at most half the gap.

Let $T=D^2\Delta^{-1}$ on $\R^4$, with its values measured in the
Frobenius norm.  Its $L^2$ operator norm is exactly one, since its
Fourier multiplier satisfies
\[
 \sum_{i,j}\left|\frac{\xi_i\xi_j}{|\xi|^2}\right|^2=1.
\]
The Calder\'on--Zygmund bound gives a finite $L^4$ operator norm
$C_4$.  Interpolation between $2$ and $4$ therefore gives
$\|T\|_{p}\le C_4^{\vartheta(p)}$, where
$1/p=(1-\vartheta(p))/2+\vartheta(p)/4$ and
$\vartheta(p)\downarrow0$ as $p\downarrow2$.  Choose once and for
all $2<p_0<4$ such that $C_4^{\vartheta(p_0)}\kappa<1$.
For a compactly supported $y$, write
$\Delta y=a:D^2y+(\Id-a):D^2y$ and absorb to obtain
\[
 \|D^2y\|_{p}\le
 \frac{\|T\|_p}{1-\kappa\|T\|_p}\|a:D^2y\|_p,
 \qquad p=2,p_0.
\]
This proves the needed Cordes estimate without differentiating $a$.

For later use its local, scaled form is
\begin{equation}\label{four:reg:cordes-local}
 \|D^2y\|_{L^p(B_{r/2})}
 \le C\left(\|a:D^2y\|_{L^p(B_r)}
                  +r^{-2}\|y\|_{L^p(B_r)}\right),
 \qquad p=2,p_0.
\end{equation}
Here is the localization.  A cutoff across a gap $h$ gives the
additional terms $Ch^{-1}\|Dy\|_p+Ch^{-2}\|y\|_p$.
On a slightly larger ball, derivative interpolation bounds the first
by $\varepsilon\|D^2y\|_p+C_\varepsilon h^{-2}\|y\|_p$.
Use concentric gaps decreasing geometrically and choose
$\varepsilon$ so that the resulting geometric series absorbs the
larger-ball Hessian norms.  For each individual function these norms
are finite, so the terminal term tends to zero.  This proves
\eqref{four:reg:cordes-local}; the same argument absorbs bounded first
order coefficients.  In particular the Christoffel terms in the
covariant equation cause no difficulty.  The analytic inputs here
are the scalar Calder\'on--Zygmund and interpolation estimates
\cite{GilbargTrudinger2001}.

\paragraph{A divergence inequality for the gradient deficit.}
Divide by $M_1+N_1$ unless this number is zero, in which case the
claim is immediate.  Subtract a constant and rescale space by $r$
and height by $r$.  The new gradient has norm at most one and the
new forcing has norm at most $r$.  Metric first derivatives and the
second fundamental form of a reflecting plane are $O(r)$;
curvature on the smooth sides is $O(r^2)$.  After normalizing the
metric at the center, all of these errors can be made uniformly
small.  In this paragraph write
$P=\grad z$, $H_z=\Hess_g z$, and $s=1-|P|_g^2\ge0$.

There is an exact cancellation of the measurable axial coefficient:
\begin{align}
 A_z\grad(|P|_g^2/2)-GP
   &=(H_z-\Delta_gz\,g)P=:Q,\label{four:reg:flux-identity}\\
 \divg Q
   &=|H_z|_g^2-(\Delta_gz)^2+\Ric_g(P,P).
                                      \label{four:reg:flux-divergence}
\end{align}
Indeed $\grad(|P|_g^2/2)=H_zP$ and
$\Delta_gz=G+(1-\widetilde\chi)H_z(e_z,e_z)$.  The axial parts
in the first formula cancel.  In the second formula,
$\divg H_z=\grad\Delta_gz+\Ric_g(P,\cdot)$ cancels the third
derivatives.  Both identities hold also at $P=0$, because the first
has both sides zero there.  At no point is a derivative of
$G$ or $\widetilde\chi$ taken.
Choose $\varepsilon_0>0$ with
$(1+\varepsilon_0)(1-\chi_0)^2<1$ when $\chi_0<1$.
Young's inequality, or directly $\Delta_gz=G$ if $\chi_0=1$, gives
\begin{equation}\label{four:reg:flux-coercivity}
 |H_z|_g^2-(\Delta_gz)^2
       \ge c|H_z|_g^2-C|G|^2,
 \qquad c=c(\chi_0)>0.
\end{equation}

We compute the only surface error in this identity.  In boundary
normal coordinates on the original side, $P=z_4\partial_4$ at the
plane and, for $a,b<4$,
\[
 (H_z)_{ab}=-\Gamma^4_{ab}z_4
             =\tfrac12(\partial_4g_{ab})z_4.
\]
Writing $H_B=\tfrac12g^{ab}\partial_4g_{ab}$ for the mean curvature
toward increasing $x_4$, this gives
\[
 Q_4=(z_{44}-\Delta_gz)z_4=-H_Bz_4^2.
\]
Thus the jump of the volume-weighted normal flux on the double is
bounded by $C|\operatorname{II}_B||P|^2$.  In particular it contains
no uncontrolled second derivative of $z$.  A bounded plane density
$j(x')$ is the distributional divergence of
$j(x')\mathbf1_{\{x_4>0\}}\partial_4$; only a normal derivative is
taken in this assertion.  After a linear coordinate change use the
corresponding constant normal direction.  Subtracting this bounded
step field handles the signed jump, without imposing any tangential
regularity on $j$.

More explicitly, put $J_g=\sqrt{\det g}$ and let
\[
 a_1^{ij}=J_g\bigl(g^{ij}-(1-\widetilde\chi)e_z^ie_z^j\bigr)
\]
be the coordinate divergence matrix. Let $j=[J_gQ^4]$ be the upper trace minus the lower trace.  Reflection
gives $Q^-=R_4Q^+$, so $j=-2J_gH_Bz_4^2$.  Then
\[
 \operatorname{div}(J_gQ)
 =J_g\bigl(|H_z|^2-(\Delta_gz)^2+\Ric_g(P,P)\bigr)
                   +j\,\dd\mathcal H^3|_{\{x_4=0\}}.
\]
Since $a_1Ds+2J_gGP=-2J_gQ$, the error field
$2J_gGP+2j(x')\mathbf1_{\{x_4>0\}}\partial_4$ cancels the
plane contribution exactly.

Consequently \eqref{four:reg:flux-identity}--\eqref{four:reg:flux-coercivity},
with volume factors included, imply on a full normalized ball
\begin{equation}\label{four:reg:deficit-inequality}
 \operatorname{div}(a_1Ds+E_1)
       \le-c_1|H_z|^2+e_1,
 \qquad
 \|E_1\|_\infty+\|e_1\|_\infty\longrightarrow0
\end{equation}
as the normalized source and geometry errors tend to zero.  The
matrix $a_1$ is bounded, symmetric and uniformly elliptic, with
constants depending only on $\chi_0$ and the fixed comparison bounds.
For example, away from the plane one may take the volume-weighted
matrix of $A_z$, $E_1=2\sqrt{\det g}\,GP$, and
$e_1\le C(|G|^2+|\Ric_g|)$; the step field just described supplies
the remaining term in $E_1$.

\paragraph{The gradient drop or affine alternative.}
Fix $r_*=1/32$.  We claim that, for every $b_0>0$, there are
$r_*<k_*<1$ and $\delta_{\rm alt}>0$, depending only on the fixed
parameters and $b_0$, such that a normalized solution on $B_1$,
with gradient at most one and errors at most $\delta_{\rm alt}$,
satisfies one of
\begin{equation}\label{four:reg:alternative}
 \sup_{B_{r_*}}|\grad z|_g\le k_*;\qquad
 \|z-L\|_{L^\infty(B_{r_*})}\le r_*b_0,
       \quad \tfrac12\le|DL|\le2,
\end{equation}
where $L$ is a coordinate affine function.

To prove the claim, consider a sequence with errors tending to zero
and $\inf_{B_{r_*}}s\to0$.  Weak Harnack for nonnegative scalar
divergence supersolutions, applied to
\eqref{four:reg:deficit-inequality} with its bounded vector and scalar
errors, yields for some $q>0$
\[
 \left(|B_{1/2}|^{-1}\int_{B_{1/2}}s^q\right)^{1/q}
 \le C\left(\inf_{B_{r_*}}s+\|E_1\|_\infty+\|e_1\|_\infty\right)
       \longrightarrow0.
\]
One obtains these fixed radii from the usual concentric weak Harnack
inequality by a finite chain of overlapping balls.  Thus $s\to0$
in measure.  The scalar weak Harnack and energy estimates used here
are valid for bounded measurable divergence matrices \cite{GilbargTrudinger2001}.

There is also $Ds\to0$ in $L^1$ on smaller balls.  Indeed, for a
fixed cutoff $\zeta$ and $b>0$, test the weak inequality by
$\zeta^2(b-s)_+$.  Discarding its favorable Hessian term and using
Young's inequality gives
\[
 \int_{\{s<b\}}\zeta^2|Ds|^2
 \le Cb^2\int|D\zeta|^2+C\|E_1\|_\infty^2
             +Cb\bigl(\|E_1\|_\infty+\|e_1\|_\infty\bigr).
\]
The full local $L^2$ norm of $Ds$ is bounded by
\eqref{four:reg:cordes-local}, the gradient bound and the geometry bounds.
On $\{s\ge b\}$, Cauchy--Schwarz and convergence in measure give
a vanishing $L^1$ norm.  On $\{s<b\}$ the displayed estimate bounds
the limiting $L^1$ norm by $Cb$.  Letting $b\downarrow0$ proves the
claim.  Testing \eqref{four:reg:deficit-inequality} with a fixed
nonnegative cutoff equal to one on $B_{1/4}$ now gives
\[
 \int_{B_{1/4}}|H_z|^2\longrightarrow0.
\]
The coordinate Hessian also tends to zero in $L^2$, because the
connection coefficients tend to zero and $Dz$ is bounded.
After subtracting $z(0)$, the functions are equi-Lipschitz.
Poincar\'e gives convergence of their gradients in $L^2$ to constants,
and compactness gives uniform convergence to an affine function on
smaller balls.  Since $s\to0$ in measure and the metrics converge to
the identity, its slope has length one.  If
\eqref{four:reg:alternative} failed for every $k_*\uparrow1$ and every
tolerance tending to zero, this conclusion would contradict failure
of its affine alternative.  This proves the claim.

\paragraph{All derivatives for a fixed exceptional axis.}
For a constant Euclidean unit vector $e_0$, let
\begin{equation}\label{four:reg:fixed-axis-operator}
 \mathcal A_0Y:=\Delta_{e_0^\perp}Y
                 +\widetilde\chi(x)\partial_{e_0e_0}Y.
\end{equation}
If $Y_0\in W^{2,2}(B_{3/4})$ and $\mathcal A_0Y_0=0$ almost
everywhere, we claim that, for some $\alpha_1\in(0,1)$,
\begin{equation}\label{four:reg:fixed-axis-estimate}
 \|DY_0\|_{C^{\alpha_1}(B_{1/4})}
       \le C\|Y_0\|_{W^{2,2}(B_{3/4})}.
\end{equation}
Set $v_1=\partial_{e_0}Y_0\in W^{1,2}$.  Differentiating the
equation in distributions, with no differentiation of its
coefficient as a separate function, gives
\[
 \Delta_{e_0^\perp}v_1+\partial_{e_0}
                (\widetilde\chi\partial_{e_0}v_1)=0.
\]
This is a scalar uniformly elliptic divergence equation.  Its local
boundedness and H\"older estimate bound $v_1$ in $C^{\alpha_1}$.
Caccioppoli, subtracting the value at the center of a ball of radius
$h$, then gives
\begin{equation}\label{four:reg:fixed-axis-morrey}
 \int_{B_h(x)}|Dv_1|^2\le Ch^{2+2\alpha_1}
                   \|Y_0\|_{W^{2,2}(B_{3/4})}^2
\end{equation}
at all centers in a fixed smaller ball.  The power follows from
$h^{-2}$ for the cutoff, $h^{2\alpha_1}$ for the squared
oscillation, and volume $h^4$.

Normalize the norm on the right to one.  Since
$\Delta Y_0=(1-\widetilde\chi)\partial_{e_0}v_1$, it follows that
\begin{equation}\label{four:reg:laplacian-morrey}
 \int_{B_h(x)}|\Delta Y_0|\le Ch^{3+\alpha_1}.
\end{equation}
Let $\mathcal N$ be the Newtonian potential of $\zeta\Delta Y_0$, where
$\zeta$ is supported inside $B_{1/2}$ and equals one on a slightly
smaller ball.  The remainder $Y_0-\mathcal N$ is harmonic there, with bounded
local $L^2$ norm; the potential has such a bound by local integrability
of $|x|^{-2}$ and the $L^2$ bound for $\Delta Y_0$.  For the gradient
of $\mathcal N$, the kernel and its derivative are bounded by $C|x|^{-3}$
and $C|x|^{-4}$.  If two points are separated by $h$, the annuli of
radii $s\le4h$ give, by \eqref{four:reg:laplacian-morrey},
$C\sum s^{\alpha_1}\le Ch^{\alpha_1}$.  On annuli $s>4h$ the
kernel difference instead gives
$Ch\sum s^{\alpha_1-1}\le Ch^{\alpha_1}$.
The sums run over dyadic scales; both converge because
$0<\alpha_1<1$.  The harmonic remainder has bounded smooth
derivatives on the final smaller ball.  Thus every component of
$DY_0$, and not only $\partial_{e_0}Y_0$, satisfies
\eqref{four:reg:fixed-axis-estimate}.

\paragraph{Improvement of a nonzero affine approximation.}
Suppose on $B_1$ that
\begin{equation}\label{four:reg:affine-input}
 \|z-L\|_\infty\le b,\quad 0<b\le b_0,\quad
 \tfrac14\le|DL|\le4,\quad |\grad z|_g\le8,
\end{equation}
and the source, metric oscillation from the identity and metric
first derivatives are at most $b\delta_{\mathrm{aff}}$.  Put $Y=(z-L)/b$.
Its coordinate equation has uniformly bounded right hand side on
$B_1$: the covariant Hessian of $L$ is a Christoffel symbol times
its bounded slope, and division by $b$ leaves a bound $C\delta_{\mathrm{aff}}$.
Consequently \eqref{four:reg:cordes-local} gives a uniform
$W^{2,p_0}(B_{7/8})$ bound for $Y$.

Let $e_0=DL/|DL|$ and let $a_0=\Id-(1-\widetilde\chi)e_0\otimes e_0$.
Since $Dz=DL+bDY$, the set
$E_b=\{|bDY|>\sqrt b\}\cap B_{3/4}$ has measure at most
$Cb^{p_0/2}$.  On its complement the actual gradient is bounded
away from zero and the coordinate principal matrix differs from
$a_0$ by at most $C(\sqrt b+b\delta_{\mathrm{aff}})$.  On $E_b$ this difference
is bounded.  H\"older's inequality, with
$1/2=1/p_0+1/q$, shows explicitly that
\begin{equation}\label{four:reg:small-fixed-source}
 \|a_0:D^2Y\|_{L^2(B_{3/4})}
 \le C\left(\delta_{\mathrm{aff}}+\sqrt{b_0}
                  +b_0^{(p_0-2)/4}\right)=:\varepsilon_1.
\end{equation}
Zeros of $Dz$ can occur only in the exceptional set when $b_0$ is
small.  The arbitrary permitted choice of their axis is therefore
covered by the same estimate.

Remove this residual by solving $a_0:D^2y_1=a_0:D^2Y$ on $B_{3/4}$
with zero Dirichlet data.  To justify this solve for measurable
$\widetilde\chi$, recall the Miranda--Talenti boundary identity on a
Euclidean ball for smooth zero-trace $y$
\cite[Section~2, Theorem~3, p.~297]{Talenti1965}:
\[
 \int\bigl((\Delta y)^2-|D^2y|^2\bigr)
       =\int_{\partial B}H_{\partial B}(\partial_\nu y)^2\ge0.
\]
Here mean curvature is the outward trace convention, equal to minus
three times Talenti's signed normalized mean curvature. The identity
extends by density to $W^{2,2}\cap W^{1,2}_0$.  Hence the
Dirichlet inverse of $\Delta$ has Hessian norm at most one.
On this space the map
\[
 y\longmapsto\Delta_D^{-1}
       \bigl(a_0:D^2Y+(\Id-a_0):D^2y\bigr)
\]
is a contraction in the Hessian $L^2$ norm with factor at most
$1-\chi_0$.  Poincar\'e and the Dirichlet Laplace estimate identify
this as a complete norm and give
\begin{equation}\label{four:reg:small-correction}
 \|y_1\|_{W^{2,2}(B_{3/4})}\le C\varepsilon_1.
\end{equation}
The function $Y_0=Y-y_1$ has uniformly bounded $W^{2,2}$ norm and
satisfies the homogeneous fixed-axis equation.  Therefore
\eqref{four:reg:fixed-axis-estimate} applies to it.

The passage to a small supremum norm must use a separate estimate
in dimension four.  Choose any
$0<\eta<\min\{1,2-4/p_0\}$.  The $W^{2,p_0}$ bound gives a
uniform $C^{0,\eta}$ modulus for $Y$ on smaller balls.  The
fixed-axis estimate and the local $L^2$ norm give a uniform
Lipschitz bound for $Y_0$.  Thus $y_1$ has a common $C^{0,\eta}$
modulus on $B_{1/4}$.  If $m=\|y_1\|_{L^\infty(B_{1/8})}$, a
ball of radius $c m^{1/\eta}$ around a point with absolute value
at least $m/2$ has absolute value at least $m/4$, provided $m$ is
small.  Hence
\[
 \|y_1\|_2^2\ge c m^{2+4/\eta},\qquad
 m\le C\varepsilon_1^{\eta/(\eta+2)}.
\]
If $m$ is not small, the same reasoning on a fixed smaller ball
first excludes that case when $\varepsilon_1$ tends to zero.
This argument uses no embedding of $W^{2,2}(\R^4)$ into $C^0$.

Choose $0<\alpha_2<\alpha_1$, then $0<\lambda_0<1/8$ so that
the Taylor remainder constant in \eqref{four:reg:fixed-axis-estimate}
satisfies $C\lambda_0^{1+\alpha_1}
\le\tfrac12\lambda_0^{1+\alpha_2}$.
Next choose $\delta_{\mathrm{aff}},b_0>0$ so that the correction supremum above
is at most $\tfrac12\lambda_0^{1+\alpha_2}$.
Taylor approximation of $Y_0$ at the center gives an affine $L_1$
such that
\begin{equation}\label{four:reg:affine-improvement}
 \|z-L_1\|_{L^\infty(B_{\lambda_0})}
       \le b\lambda_0^{1+\alpha_2},\qquad
 |DL_1-DL|\le Cb.
\end{equation}
Reduce $b_0$ further so that
\begin{equation}\label{four:reg:slope-sum}
 Cb_0/(1-\lambda_0^{\alpha_2})<1/4.
\end{equation}
These estimates concern the almost everywhere coordinate equation
on full balls and therefore apply equally to doubled patches.

\paragraph{The complete scaling iteration.}
The choices are ordered as follows: first $p_0,\alpha_1$ and their
constants; then $\alpha_2,\lambda_0,\delta_{\mathrm{aff}},b_0$; then
$k_*,\delta_{\rm alt}$ from \eqref{four:reg:alternative} for that $b_0$.
Finally shrink the initial spatial scale so that its normalized
source and geometry errors are bounded by $\delta$, where
\begin{equation}\label{four:reg:initial-tolerance}
 \delta\le c\min\{\delta_{\rm alt},b_0\delta_{\mathrm{aff}}/r_*\}.
\end{equation}
Here the fixed $c>0$ incorporates all comparison constants in the
geometry rescalings.  This scale depends only on the theorem's
parameters, since division by $M_1+N_1$ makes the initial source
at most one.  Write the resulting function as $u$ and center all
following balls at zero.

After $m$ successive gradient drops the actual rescaled solution is
\[
 u_m(x)=\frac{u(r_*^m x)-u(0)}{r_*^m k_*^m},\qquad
 |\grad u_m|_{g_m}\le1,\qquad
 \|G_m\|_\infty\le\delta(r_*/k_*)^m.
\]
The metric remains the identity at the center; its oscillation and
first derivative errors are at most $C\delta r_*^m$, with its
curvature error shrinking at least as fast.  A reflecting plane
has the correspondingly scaled second fundamental form.  Since
$r_*<k_*$, the next application of the alternative is legitimate.
If the affine alternative first occurs after these $m$ drops,
rescale its ball by setting $v_0(x)=u_m(r_*x)/r_*$.
Then $v_0$ has affine error at most $b_0$, slope in $[1/2,2]$,
bounded gradient, and source at most
\[
 r_*\delta(r_*/k_*)^m\le b_0\delta_{\mathrm{aff}}.
\]
Its geometry satisfies the same tolerance by
\eqref{four:reg:initial-tolerance}.

In all subsequent affine steps use
\[
 v_i(x)=\lambda_0^{-i}v_0(\lambda_0^i x),\qquad
 b_i=b_0\lambda_0^{i\alpha_2}.
\]
The derivative $Dv_i(x)=Dv_0(\lambda_0^i x)$ retains the actual
gradient bound.  The source and first order geometry errors are
at most $b_0\delta_{\mathrm{aff}}\lambda_0^i\le b_i\delta_{\mathrm{aff}}$, because
$\alpha_2<1$.  The affine function used at stage $i$ is rescaled
in the same way, including its constant term.  Thus
\eqref{four:reg:affine-improvement} applies inductively.  Its changes
of slope sum to less than $1/4$ by \eqref{four:reg:slope-sum}, so all
slopes remain in $[1/4,4]$.  Only the auxiliary residual in the
proof of an improvement is divided by $b_i$; that residual is never
asserted to solve the original nonlinear axial equation.

Choose
\[
 0<\alpha\le\min\{\alpha_2,\log k_*/\log r_*\}.
\]
If the drops never end, the gradient on $B_{r_*^m}$ is at most
$k_*^m\le r_*^{m\alpha}$, and subtraction of the central value
gives a constant affine approximation with error $Cr_*^{m(1+\alpha)}$.
If the drops end after $m$ steps, the affine approximation at
$R_i=r_*^{m+1}\lambda_0^i$ has error at most
\[
 b_0 k_*^m r_*^{m+1}\lambda_0^{i(1+\alpha_2)}
       \le b_0r_*^{-\alpha}R_i^{1+\alpha}.
\]
The previous drop scales have the same type of estimate.  Adjacent
radii have ratio $r_*$ or $\lambda_0$, both fixed, so every
intermediate radius has an affine approximation with error
$Cr^{1+\alpha}$.  This conclusion holds at every center of the
smaller full patches, including centers whose balls cross the plane.

For clarity, the usual comparison of affine approximations gives
the gradient conclusion directly.  If $L_r,L_{r/2}$ are two such
approximations, then
$\|L_r-L_{r/2}\|_{L^\infty(B_{r/2})}\le Cr^{1+\alpha}$, whence
$|DL_r-DL_{r/2}|\le Cr^\alpha$.  Their slopes converge, and
the limit is $Dz$ since $z$ is $C^1$.  At points separated by $h$,
compare the two approximations on overlapping balls of radius $4h$;
the same affine norm estimate gives a slope difference $Ch^\alpha$.
Their limits differ from those slopes by $Ch^\alpha$ as well.
This proves the first estimate in \eqref{four:eq:axial-estimate}.

Finally take an affine approximation $L$ on $B_r$ with error
$Cr^{1+\alpha}$.  Its $L^2$ error is at most $Cr^{3+\alpha}$.
The coordinate equation for $z-L$ has right hand side equal to
$G$ plus a bounded Christoffel coefficient times $Dz$, whose
$L^2$ norm is at most $Cr^2$ before these iterative rescalings.
Equation \eqref{four:reg:cordes-local} with $p=2$ therefore gives
\[
 \|D^2z\|_{L^2(B_{r/2})}\le C(r^{1+\alpha}+r^2).
\]
The difference $\Hess_gz-D^2z$ has $L^2$ norm at most $Cr^2$.
Since $\alpha<1$, squaring and restoring $M_1+N_1$ proves the
second estimate in \eqref{four:eq:axial-estimate}.
\end{proof}

\subsection{The measure estimate needed in dimension four}

The second regularity step is Lemma~\ref{b:lem:natural-growth} with
$n=4$. Its input is a bounded scalar function satisfying
\begin{equation}\label{four:reg:natural-equation}
 \operatorname{div}(\mathcal A DZ+B_1)=E,
\end{equation}
with symmetric measurable uniformly elliptic $\mathcal A$, bounded
$B_1$, and a signed source dominated by
\begin{equation}\label{four:reg:natural-growth-bound}
 |E|\le C_0(1+|DZ|^2)\dd x+\mu_1,\qquad
 \mu_1(B_r(x))\le C_\mu r^{2+\beta},\qquad \beta>0.
\end{equation}
The preceding axial theorem will supply the Hessian part of $\mu_1$.
A bounded density on a reflecting three-plane contributes $O(r^3)$,
which also has the required positive exponent. The shared lemma
absorbs $|DZ|^2$ by its exponential product rule and bounds the
remaining Green correction by $C(r+r^\beta)$.

In the application below each individual solution is smooth, and its
reflection is continuous and smooth on the two closed sides. Its common
trace therefore defines the measure products and verifies that product
rule, including the plane flux. The H\"older representative supplied
by the lemma agrees everywhere with this continuous function.
We now verify the coefficients, the reflected source, and the ball-growth
bound for the actual homotopy before applying the estimate.

\subsection{Application to the homotopy and the boundary estimates}

\begin{proposition}\label{four:prop:coupled-regularity}
Fix the prepared data, $\eps$, a sufficiently large $N$ and any
stage of the homotopy.  Every smooth solution on an allowed finite
truncation $\Omega_R$ with $t\ge-\eps$ has uniform local estimates
for every finite number of derivatives of $f$ and $Z$, up to both
the inner and outer boundary faces.  The constants on the fixed
compact part and on uniform unit patches of the end are independent
of $R$ and of the homotopy parameters.  They may depend arbitrarily
on $N$.  On each fixed finite truncation they therefore bound
$\|Z\|_{C^{1,\alpha_0}}$ for every prescribed $0<\alpha_0<1$.
\end{proposition}

\begin{proof}
We use the uniform family of interior and boundary patches described
in the statement, shrinking them by fixed factors when necessary.

\paragraph{The first H\"older estimates and the reflected equation.}
The bounded range \eqref{four:eq:bounded-range} keeps
$f,Z,t,\grad f$ bounded and $l,d,\chi,u$ above positive constants
depending on $N$.  Divide the trace equation by $l/\sqrt D$ to obtain
\begin{equation}\label{four:reg:trace-normalized}
 \begin{gathered}
 \Achi:\Hess_g f=G_f,\qquad
 G_f=\frac{\sqrt D}{l}
          \bigl(F_{\rm presc}-\tr_{\Achi}K\bigr),\\
 |G_f|\le C(N),\qquad 0<\chi_0(N)\le\chi\le1.
 \end{gathered}
\end{equation}
All prescribed trace expressions have bounded values on this range.
Every boundary value of $f$ is constant on its face.  Thus
Theorem~\ref{four:thm:axial-regularity} gives a H\"older bound for $Df$
and
\begin{equation}\label{four:reg:f-morrey}
 \int_{B_r(x)\cap\Omega_R}|D_x^2f|^2\,\dd x
                         \le C(N)r^{2+2\alpha}.
\end{equation}
The coordinate and covariant Hessians differ by a bounded
Christoffel term times $Df$, so they have the same bound at these
radii.  The bound also holds for the odd double at a face.

In coordinates set
\[
 \mathcal A^{ij}=3\sqrt{\det g}\,u\,A_\chi^{ij},\qquad
 B_1^i=\lambda\sqrt{\det g}\,u\,
                     \bigl(\Achi K(w,\cdot)\bigr)^i,
\]
using $\lambda=0$ in the later stages.  The coordinate $Z$ equation
is \eqref{four:reg:natural-equation}, with these bounded coefficients
and the appropriate scalar homotopy factor in $E$.  Its matrix
is symmetric and uniformly elliptic at fixed $N$.
Differentiating the implicit relation for $t$ gives the exact
covector identity
\begin{equation}\label{four:reg:t-differential}
 3\,\dd Z=(3+pv)\,\dd t+H^f(w,\cdot).
\end{equation}
Since $3+pv\ge3$ and all its coefficients are bounded on the range,
$|D_xt|\le C(N)(1+|D_xZ|+|D_x^2f|)$ in coordinates.
The quadratic formula \eqref{four:eq:quadratic-form} and the bounded
lower order terms in the source therefore give
\begin{equation}\label{four:reg:coupled-source}
 |E|\le C(N)(1+|D_xZ|^2+|D_x^2f|^2)
\end{equation}
in the interior.  No derivative of the principal matrix has been
used to obtain this bound.

At a face take the domain to be $x_4>0$ and write
$Q=\mathcal A DZ+B_1$.  For $x_4<0$ define, with $s=1$ or $-1$,
\[
 \widetilde Z(x)=sZ(R_4x),\quad
 \widetilde{\mathcal A}(x)=R_4\mathcal A(R_4x)R_4,\quad
 \widetilde B_1(x)=sR_4B_1(R_4x).
\]
The reflected flux is $\widetilde Q(x)=sR_4Q(R_4x)$.
Integration by parts on the two sides shows that its divergence
has the reflected bulk source $sE(R_4x)$ and the plane term
\begin{equation}\label{four:reg:reflection-measure}
 (1+s)Q_4^+\,\dd\mathcal H^3|_{\{x_4=0\}}.
\end{equation}
Indeed the upper normal component is $Q_4^+$ and the lower one is
$-sQ_4^+$, so their difference is $(1+s)Q_4^+$.
For an inner face use even reflection, $s=1$.  The total conormal
flux is prescribed by \eqref{four:eq:homotopy-flux}; its coordinate
density $Q_4^+$ is bounded on the established range.  For an outer
face use odd reflection, $s=-1$.  The zero Dirichlet condition
makes $Z$ continuous across the face, and the plane term vanishes.
Even reflection at the inner face is also continuous.  Each
individual reflected function belongs to $W^{1,2}$ and is smooth
on the two closed sides, which verifies the weak chain rule
required in Lemma~\ref{b:lem:natural-growth}.

In the full patch, \eqref{four:reg:coupled-source} and
\eqref{four:reg:reflection-measure} give a signed measure dominated by
\[
 C(N)(1+|D_xZ|^2)\,\dd x+\mu_1,\qquad
 \mu_1=C(N)|D_x^2f|^2\,\dd x
               +2|Q_4^+|\,\dd\mathcal H^3|_{\{x_4=0\}},
\]
where the Hessian density is reflected and the plane term is
omitted for an odd face or an interior patch.  By
\eqref{four:reg:f-morrey}, this positive measure satisfies, at every
center,
\[
 \mu_1(B_h(x))\le C(N)h^{2+\beta},\qquad
                   \beta=\min\{2\alpha,1\}>0.
\]
For balls crossing the plane this follows either from the doubled
estimate or by enlarging a ball centered on the plane by a fixed
factor.  Lemma~\ref{b:lem:natural-growth} now supplies a uniform
H\"older modulus for $Z$, including at both boundary types.

\paragraph{Schauder for $f$ and the ordinary normal condition for $Z$.}
The actual coefficients in \eqref{four:reg:trace-normalized} are smooth
functions of $(x,Z,Df)$ on the bounded range.  This remains true
at $Df=0$, as is manifest from
\[
 \Achi=\Id-\frac{3+p}{3+pv}w\otimes w.
\]
The source is smooth in these variables and $f$, uniformly over
the four parameter ranges.  The H\"older bounds just proved
therefore make the coefficients and source uniformly
$C^{0,\theta}$ for some $\theta>0$.  The scalar Dirichlet Schauder
estimate, with the constant face data, gives
\begin{equation}\label{four:reg:f-schauder}
 \|f\|_{C^{2,\theta}(U')}\le C(N)
\end{equation}
on every smaller interior or boundary patch \cite{GilbargTrudinger2001}.
This application has H\"older principal coefficients, a H\"older
right hand side, smooth boundary, and uniform ellipticity; it
does not apply a system estimate.

Now expand the divergence equation for $Z$.  Derivatives of its
coefficients, which depend smoothly on $(x,Z,Df)$, are sums of
bounded terms and bounded multiples of $DZ,D^2f$.  The latter
Hessian is bounded by \eqref{four:reg:f-schauder}.  Thus
\begin{equation}\label{four:reg:Z-nondivergence}
 a_2^{ij}(x)D_{ij}Z=R_2,\qquad
 |R_2|\le C(N)(1+|DZ|^2),
\end{equation}
with bounded $C^{0,\theta}$ uniformly elliptic principal
coefficients.  The source contains no second derivative of $Z$.

On an inner face $Df$ is normal and hence $\Achi\nu=\chi\nu$.
Dividing the prescribed flux by $3u\chi>0$ gives
\begin{equation}\label{four:reg:normal-data}
 \partial_\nu Z=b_2(x,Z,f,Df).
\end{equation}
The function $b_2$ is smooth on the bounded range, separately on
each homotopy stage, with uniform bounds through their matching
endpoints.  Extend its displayed expression smoothly in the
spatial variable into the collar.  Equation
\eqref{four:reg:f-schauder} then gives
\[
 |D_xb_2(x,Z,f,Df)|\le C(N)(1+|D_xZ|).
\]
Consequently the Sobolev trace theorem bounds the
$W^{1-1/p,p}$ norm of this normal datum by
$C(N)(1+\|Z\|_{W^{1,p}})$ on a larger patch, for each finite
$p>1$.  This derivative count costs $DZ$ and $D^2f$, never $D^2Z$.

\paragraph{Linear estimates and absorption of the quadratic term.}
Fix $p>4$.  On a ball or smooth half-patch of size $h$, the local
linear $W^{2,p}$ estimate for \eqref{four:reg:Z-nondivergence}, with
zero outer Dirichlet data or \eqref{four:reg:normal-data}, has the form
\begin{equation}\label{four:reg:Z-linear}
 \|D^2Z\|_{L^p(Q_h)}
 \le C_*\|DZ\|_{L^{2p}(Q_{2h})}^2
          +C_h\bigl(1+\|Z\|_{W^{1,p}(Q_{2h})}\bigr).
\end{equation}
Here $C_*$ is independent of sufficiently small $h$; $C_h$ may
contain inverse powers of $h$.  The constants are uniform over
the locations and the truncations in question.

We specify the linear estimate underlying this assertion.  Freeze
the continuous principal matrix at the patch center.  The interior
constant-coefficient Hessian estimate follows from the Euclidean
one by a fixed linear change of coordinates.  At the boundary,
normal coordinates give no mixed normal--tangential entries in
the frozen matrix, because its exceptional axis is normal there.
A tangential linear transformation and a normal dilation reduce
the frozen principal part to the Laplacian and preserve the plane.
Subtract a Sobolev extension of the normal datum, or of the
Dirichlet datum, to obtain homogeneous boundary data.  Even
reflection for homogeneous normal data and odd reflection for
homogeneous Dirichlet data give the constant-coefficient half-space
estimate.  The extension costs precisely the
$W^{1-1/p,p}$ norm for the normal datum.  Coefficient oscillation
contributes $\sup|a_2-a_2(x_0)|\,\|D^2Z\|_p$ and is absorbed on
small patches, uniformly by the established H\"older bound.
Cutoffs and interpolation handle the first order terms.  Scaling
shows that the coefficient of the interior $L^p$ source norm is
the constant $C_*$, independent of $h$; the lower order and trace
terms account for $C_h$.  This proves
\eqref{four:reg:Z-linear} by the scalar local estimates
\cite{GilbargTrudinger2001}.  The construction requires no meeting of different
boundary types, and the inner and outer faces here are disjoint.

For use in its nonlinear right hand side, the scaled scalar
interpolation inequality on $Q_h$ is
\begin{equation}\label{four:reg:quadratic-interpolation}
 \|DZ\|_{L^{2p}(Q_h)}^2
 \le C\,\operatorname{osc}_{Q_h}Z\,
                  \|D^2Z\|_{L^p(Q_h)}
       +Ch^{4/p-2}\bigl(\operatorname{osc}_{Q_h}Z\bigr)^2.
\end{equation}
One can obtain it on the unit patch by subtracting a constant,
using a bounded Sobolev extension on a smooth ball or half-patch,
and applying the bounded-domain Gagliardo--Nirenberg inequality
with derivative orders $1,2$, exponents $2p,p,\infty$ and parameter
$1/2$. A convenient modern formulation is
\cite[Theorem~1.3]{LiZhang2022}.  To see the interpolation mechanism directly
for a compactly supported function $v$, integration by parts gives
\[
 \int|Dv|^{2p}
 \le C_p\|v\|_\infty
             \int|Dv|^{2p-2}|D^2v|
 \le C_p\|v\|_\infty
      \left(\int|Dv|^{2p}\right)^{1-1/p}\|D^2v\|_p.
\]
Division gives the required product estimate.  Extension and
cutoff produce the additional $\|v\|_\infty^2$ term on a bounded
unit patch. This lower-order term is essential: the bounded-domain
Gagliardo--Nirenberg estimate includes an additive multiple of
$\|v\|_\infty$, so one may not drop its square after applying the
estimate to $v=Z-c$. Taking $c$ between the maximum and minimum replaces
this norm by the oscillation used here.  Finally rescaling $x=hy$ multiplies a squared
$L^{2p}$ gradient norm and the Hessian product by $h^{4/p-2}$,
which proves precisely the last power in
\eqref{four:reg:quadratic-interpolation}.

Here is a cover argument that absorbs larger-patch norms without
assuming their uniform boundedness in advance.  Use uniformly
comparable patches of radius $h$ covering the finite truncation,
and their fixed-factor enlargements.  Each enlargement is covered
by at most $M_{\mathrm{cov}}$ smaller patches, where $M_{\mathrm{cov}}$ is independent of $h$,
$R$ and the total number of patches.  Let
\[
 S=\sup_j\|D^2Z\|_{L^p(Q_h^j)}.
\]
For each individual smooth solution on the finite truncation,
$S$ is finite.  The H\"older estimate for $Z$ gives
$\operatorname{osc}_{Q_{2h}^j}Z\le C(N)h^\gamma$.
Apply \eqref{four:reg:quadratic-interpolation} in the enlarged patches
and then cover them by smaller ones.  Its contribution to
\eqref{four:reg:Z-linear} is at most
$C_*C(N)h^\gamma M_{\mathrm{cov}}^{1/p}S+C(N,h)$.
Choose $h$ small enough, also satisfying the linear freezing
condition, that this coefficient is at most $1/4$.
For the remaining first order term, local interpolation gives,
for any $\varepsilon>0$,
\[
 \|DZ\|_{L^p(Q_{2h}^j)}
 \le\varepsilon\|D^2Z\|_{L^p(Q_{2h}^j)}
                   +C(h,\varepsilon)\|Z\|_{L^p(Q_{2h}^j)}.
\]
Choose $\varepsilon$ so that $C_h\varepsilon M_{\mathrm{cov}}^{1/p}\le1/4$.
The bounded range controls the last term.  Taking the supremum
in \eqref{four:reg:Z-linear} now gives
$S\le C(N,h)+S/2$, and hence a uniform bound for $S$.
First order interpolation supplies the corresponding local
$W^{2,p}$ bounds for $Z$.  All constants may depend on the now
fixed $h$ and on $N$, but not on the total number of patches or $R$.

\paragraph{Alternating Schauder estimates and uniform geometry.}
Since $p>4$, Sobolev embedding gives $Z\in C^{1,\theta'}$ for
some $\theta'>0$.  In conjunction with \eqref{four:reg:f-schauder},
the trace equation now has $C^{1,\theta''}$ coefficients and source
for $0<\theta''\le\min\{\theta,\theta'\}$.  One differentiation
and the scalar Dirichlet Schauder estimate give
$f\in C^{3,\theta''}$.  The equation for $Z$ then has
$C^{0,\theta''}$ source, and its normal datum
\eqref{four:reg:normal-data} is $C^{1,\theta''}$.
The scalar normal-derivative or Dirichlet Schauder estimate gives
$Z\in C^{2,\theta''}$.  The boundary estimate is justified by the
same frozen half-space problem used above: its normal derivative
does not annihilate a nonzero decaying solution of the homogeneous
principal equation.  Smooth normal coordinates and smooth normal
field therefore satisfy its obliqueness condition.

Repeating in the order $f$ then $Z$ increases the number of
derivatives by one each time.  The trace equation uses $Z$ without
derivatives; the expanded second equation uses only $D^2f$ and
$DZ$ in its lower terms; and its boundary expression uses only
$Df$.  These are exactly the derivative counts needed for the
iteration.  The smooth fixed coefficients give all finite orders.

The compact inner faces have uniform smooth collars.  On the
prepared end, the metric and all fixed coefficient fields have
uniform bounds on unit patches.  Large coordinate spheres have
boundary normal charts with uniform bounds of each finite order,
and are separated from the inner faces.  Thus every radius,
ellipticity bound, chart constant and local covering number used
above is independent of the sufficiently large truncation radius.
Finally a uniform local $C^2$ bound controls any prescribed
$C^{1,\alpha_0}$ norm, $\alpha_0<1$, on a fixed finite truncation.
This proves the proposition.
\end{proof}

\begin{remark}\label{four:reg:scalar-trial-interface}
Theorem~\ref{four:thm:axial-regularity} also applies in the scalar trial
problem below. Its separate height and gradient bounds give bounded
source and a positive axial eigenvalue without the floor. Interpolation
of the principal matrix with $\Id$ preserves its three unit transverse
eigenvalues. A trial $Z\in C^{1,\alpha_0}$ then supplies the coefficient
regularity for scalar Schauder estimates through $f\in C^{3,\alpha_0}$;
this requires $DZ$, but no second derivative of the trial input.
\end{remark}
